This thesis is an original work by Trung Quan Nguyen, carried out under the supervision of Professor Linglong Kong and Professor Bei Jiang. The author formulated the research questions, designed and carried out the analysis and experiments, and wrote the text; the supervisors provided guidance and feedback throughout.

The thesis is based on the manuscript ``How Many Demonstrations Is a Task Description Worth? Specificity and Reliability in In-Context Learning'' by Bruce Nguyen, the name under which the author writes, prepared in the format of the International Conference on Machine Learning. Chapters~\ref{ch:intro} and~\ref{ch:setup} to~\ref{ch:conclusion} and Appendix~\ref{app:llm-tables} follow that manuscript's text, with its appendices placed in the chapters they support and some added detail; Sections~\ref{sec:outline} and~\ref{sec:future} are new. Chapter~\ref{ch:background} is new, except that it incorporates the manuscript's review of related work, placed in the sections whose topics it concerns. The manuscript has not yet been published.

The research involved no human participants, animals, or personal data, so no research ethics approval was required.

Beyond studying large language models in Chapter~\ref{ch:llm}, the author used tools built on them, Anthropic's Claude Code and OpenAI's Codex, to help write and run experiment code, check derivations, and draft and edit text. The author designed the study, reviewed all code, results, and text, and takes responsibility for them.

The author's salary and the Claude Code and Codex accounts used in this research were paid for by Professor Andrew McGee of the Department of Economics, from funding for his own research.

The author founded Decaude, and this research is part of its research program; the manuscript lists the author's affiliations as Decaude and the University of Alberta.
